% ==============================================================================
\section{Solving Approaches}
\label{sec:approaches}
% ==============================================================================

The MILP of \S\ref{sec:milp} is exact but its size grows quickly with the
number of aircraft and blocking arcs.  We therefore develop three
complementary methods---a topology-aware GRASP, a fixed-assignment
scheduler, and a safe pipeline composition---that operate at different
abstraction levels: the GRASP works on the position assignment, the
fixed-assignment scheduler refines the timing given an assignment, and the
safe pipeline combines them.

% ------------------------------------------------------------------------------
\subsection{Topology-Aware GRASP}
\label{sec:topo_heuristic}
% ------------------------------------------------------------------------------

The driving observation behind this method is that assigning a heavy
aircraft (large $D_r$) to a position with high outgoing blocking degree
forces its long occupancy window to overlap with aircraft serviced behind
it, triggering movements.  We exploit the \emph{blocking load} $\ell(p)$ introduced in
\S\ref{sec:problem} --- the number of positions that lie downstream of
$p$ in the blocking-arc graph $G$.  The construction step penalises each
candidate pair $(r, p)$ proportionally to $D_r \cdot \ell(p)$, steering
heavy aircraft towards positions with low (ideally zero) blocking load.

Given the cost function
\[
  c(r, p) = W^D \hat{v}^D_r + W^M \hat{m}
            + W^S \hat{n}_{rp} + W^T D_r \ell(p),
\]
in which hatted quantities are estimated by a fast greedy scheduler and
$W^T$ is the topology penalty weight, the heuristic constructs a feasible
solution by biased-random greedy insertion, then improves it through a
local-search portfolio.  Diversification is obtained by restarting from
independent random seeds.  Algorithm~\ref{alg:topo} states the full
procedure.

\begin{algorithm}[htbp]
\DontPrintSemicolon
\KwIn{Instance $I$, weights $(W^M, W^D, W^S, W^T)$,
      RCL parameter $\alpha$, time budget $T$, number of starts $n_s$.}
\KwOut{Best solution $z^*$ found.}
$z^* \leftarrow +\infty$; $T_s \leftarrow T / n_s$\;
\For{$k = 1, \ldots, n_s$}{
  $\mathcal{U} \leftarrow \calR$ sorted by $D_r$ descending\;
  \While{$\mathcal{U} \neq \emptyset$}{
    evaluate $c(r, p)$ for all $(r, p) \in \mathcal{U} \times \calP_{\text{free}}$\;
    $\text{RCL} \leftarrow \{(r, p) : c(r, p) \le c_{\min} + \alpha (c_{\max} - c_{\min})\}$\;
    pick $(r, p)$ uniformly from $\text{RCL}$ and commit it; remove $r$ from $\mathcal{U}$\;
  }
  $z \leftarrow$ schedule with current assignment\;
  \Repeat{no operator improves $z$ \emph{or} slice $T_s$ exhausted}{
    apply best of: single-move, 2-opt swap, intra-position reorder,
                   perturb-and-rebuild\;
  }
  \If{$z < z^*$}{$z^* \leftarrow z$}
}
\Return $z^*$\;
\caption{Topology-aware GRASP with multi-start ($n_s$ restarts).}
\label{alg:topo}
\end{algorithm}

All four operators share the same skeleton: propose a candidate
modification of the current incumbent, rebuild the schedule from
scratch with the modified assignment (or ordering), evaluate the
objective \eqref{eq:milpobj}, and accept the first strictly improving
proposal.  On acceptance the search restarts from the cheapest
operator, so cheap moves get repeated chances to fire before the
expensive ones are tried again.  They differ in the perturbation they
apply.

The \emph{single-move} operator iterates over every aircraft and, for
each, over every alternative position, producing $O(R\,(P-1))$
candidate moves.  Each candidate reassigns one aircraft and leaves the
rest unchanged.  This is the cheapest perturbation, so it is tried
first in every iteration.

The \emph{2-opt swap} operator is a position-exchange move that
preserves the multiset of positions occupied.  It enumerates every
ordered pair $(r, r')$ of aircraft assigned to \emph{different}
positions $p$ and $p'$ and produces a candidate in which $r$ takes
position $p'$ and $r'$ takes position $p$, with all other aircraft
fixed.  Unlike single-move, the swap conserves the load distribution
across positions and therefore reaches solutions that no chain of
single-moves can reach without first overloading a position; this
proves particularly useful on the \emph{Two rows} and \emph{Triangle}
topologies, where balanced loads are critical.  The enumeration is
$O(R^2)$ candidate pairs.

The \emph{intra-position reorder} operator does not change any
position assignment.  For every position $p$ with at least two
assigned aircraft, it generates all adjacent transpositions of the
service order at $p$, plus all single-slot insertions and all
non-adjacent transpositions within the same position; it also tries
three structured reorderings per position (earliest-due-date,
slack-ascending, and delay-ratio-descending).  Because aircraft
durations and time windows interact with the precedence chain, two
schedules that differ only in the service order at one position
can produce substantially different delays and movements, so this
family of moves is essential for tightening the timing once the
assignment is good.

Finally, the \emph{perturb-and-rebuild} operator is an aggressive
delay-driven move designed to escape local minima where single-move
and 2-opt are stuck.  It identifies the
$K = \max(2, \lceil R/10 \rceil)$ aircraft with the largest current
delay, and for each of them tries (i) every intra-position slot
relocation at its current position, and (ii) every single-move
relocation to a different position.  Effectively this is a focused
combination of intra-position insertion and single-move applied only
to the worst-performing aircraft, where it has the greatest chance of
finding improvement.  The intra-position part addresses the case
where the aircraft is simply scheduled too late at its current
position; the cross-position part covers the case where the position
itself is wrong.  The move is more expensive than single-move per
candidate but evaluates only $K \cdot (P + k_p)$ candidates per pass,
which on benchmark instances ($K \le 3$, $P = 5$) is comparable in
cost to one single-move pass.

All operators dominate the construction phase in cost only because
the schedule rebuild is itself non-trivial; the move-enumeration
budget is $O(R^2)$ per pass at worst, and the local-search phase
typically converges within two to four passes before the slice
budget $T_s = T/n_s$ is exhausted.

% ------------------------------------------------------------------------------
\subsection{Fixed-Assignment Scheduler (FAS)}
\label{sec:fas}
% ------------------------------------------------------------------------------

Given a fixed position assignment $\pi : \calR \to \calP$, the AHPP reduces
to a pure scheduling problem: find aircraft start times $\{s_r\}$ that
minimise the objective \eqref{eq:milpobj} subject to the constraints
\eqref{eq:dur}--\eqref{eq:uOut}, with every $x_{rp}$ replaced by the
constant $\mathbf{1}[\pi(r) = p]$.  Two simplifications follow.  The
$O(RP)$ assignment binaries disappear, and the position-guard factors in
the blocking constraints collapse to known scalars; the blocking-pair
enumeration therefore shrinks from $O(R^2 |\calA| P^2)$ to
$O(R^2 |\calA|)$.  Job-level timings are recovered from $s_r$ by walking
the precedence chain, exactly as in \S\ref{sec:job_recovery}.  The
resulting model is encoded directly in \texttt{gurobipy}, bypassing the
Pyomo layer.

Algorithm~\ref{alg:fas} summarises the procedure.  A multi-candidate
extension allows a small set of positions $\kappa(r) \supseteq \{\pi(r)\}$
for each aircraft, recovering a fraction of the full assignment problem
while keeping the model tractable.  In that variant the candidate set is
non-trivial only for the four aircraft with the highest score
$\sigma(r) = v^D_r + 0.5 \cdot \mathbf{1}[f_r \approx m]$, and is built as
the union of their primary position and the positions they occupy in
ten diverse topology assignments produced by 1-second GRASP runs.

\begin{algorithm}[htbp]
\DontPrintSemicolon
\KwIn{Instance $I$, assignment $\pi$ (optionally candidate sets $\kappa$),
      time budget $T$.}
\KwOut{Best feasible schedule found within $T$.}
\If{candidates $\kappa$ provided}{$\mathcal{Y} \leftarrow \{x_{rp} : p \in \kappa(r)\}$}
\Else{$\mathcal{Y} \leftarrow \emptyset$ \emph{(all $x_{rp}$ fixed to $\pi$)}}
build reduced model with aggregated blocks $[s_r, s_r + D_r]$
       and variables $\mathcal{Y}$\;
solve with Gurobi under time limit $T$\;
recover job start times $t^S_j$ from $s_r$ along each precedence chain\;
\Return resulting schedule\;
\caption{Fixed-Assignment Scheduler (FAS) with optional multi-candidate.}
\label{alg:fas}
\end{algorithm}

% ------------------------------------------------------------------------------
\subsection{Safe Pipeline}
\label{sec:safe_pipeline}
% ------------------------------------------------------------------------------

The two preceding methods can be run in sequence: the GRASP produces a
position assignment, and FAS refines its timing.  Their relative quality
is instance-dependent, since neither dominates the other for every
configuration.  The \emph{safe pipeline} therefore returns the best
objective achieved by any of the component methods:
\[
  z^* = \min_{h \in \mathcal{H}} z_h,
  \qquad \mathcal{H} = \{\texttt{topology\_ms6},\ \texttt{fas\_on\_topo}\},
\]
optionally extended with \texttt{fas\_2cand}.  Since each component
already stores its solution in a shared cache, the pipeline performs no
additional optimisation; its cost is a single comparison.

By construction, the safe pipeline objective satisfies $z^* \le z_h$ for
every $h \in \mathcal{H}$; in particular
$z^* \le z_{\texttt{topology\_ms6}}$, so the pipeline never degrades the
GRASP result.
